\clearpage\fancyhead[L]{\small COMET.}\hypertarget{native-350-lower}{}\label{native-350-lower}
\begin{center}\textbf{COMET.}\par\vspace{4pt}\hrule\vspace{4pt}\fontsize{10}{12}\selectfont Göttingische gelehrte Anzeigen. No. 15, pp. 147..149, 25 January 1812.\par\vspace{4pt}\hrule\vspace{4pt}\end{center}

\par\indent The weather has been unusually unfavorable for more than a month. Since the first observations of the \emph{new} comet reported in these pages (no.~201 of last year [p.~342 of this volume]), it has allowed only three further observations at this observatory. We give them here together with the first observations, whose reduction has been slightly improved.

\begin{center}
\fontsize{10}{12}\selectfont\setlength{\tabcolsep}{3pt}\renewcommand{\arraystretch}{1.15}\begin{tabular}{l|c|c|c}
 & \shortstack{Mean\\time} & \shortstack{Right\\ascension} & Declination \\
\hline
1811 Dec. 9 & $10^{\text{u}}6^{\text{m}}52^{\text{s}}$ & $63^\circ49'41''4$ & $10^\circ21'55''5$ S. \\
11 & $10\;\;34\;\;1$ & $63\;\;33\;\;18.0$ & $8\;\;39\;\;46.4$ \\
12 & $8\;\;5\;\;52$ & $63\;\;26\;\;25.3$ & $7\;\;54\;\;25.9$ \\
1812 January 3 & $7\;\;1\;\;30$ & $62\;\;37\;\;6.4$ & $9\;\;47\;\;14.7$ N. \\
4 & $6\;\;53\;\;59$ & $62\;\;40\;\;38.6$ & $10\;\;29\;\;57.8$ \\
\end{tabular}
\end{center}

\noindent As early as 13 December, Professor \textsc{Gauss} had calculated provisional parabolic elements from his first observations and from those of Mr. von \textsc{Zach}, which we also reported in the same place. These elements appeared in the Westphälischer Moniteur, no.~301, and in the December issue of the Monatliche Correspondenz. Naturally, elements based on such scant data could only be approximate, but on 3 January they still differed from the observation by only 8 minutes. After making his observation of 4 January, Professor \textsc{Gauss} entrusted the refinement of the parabolic elements to Mr. \textsc{Nicolai}, who carried out the work with the greatest care and with the best results. The\SourceLeaf{351} following are the improved elements this skilled calculator obtained. They will surely admit only minor further improvements.

\Needspace{15\baselineskip}
\begin{center}\fontsize{10}{12}\selectfont\setlength{\tabcolsep}{4pt}\renewcommand{\arraystretch}{1.15}\begin{tabular}{l}
Time of perihelion 1811 Nov. 11\\
\quad $9^{\mathrm u}\;24^{\mathrm m}\;38^{\mathrm s}$ Göttingen meridian.\\
\end{tabular}\end{center}
\begin{center}\fontsize{10}{12}\selectfont\setlength{\tabcolsep}{3pt}\renewcommand{\arraystretch}{1.15}\begin{tabular}{llr}
Perihelion distance & .\;.\;.\;.\;.\;.\;.\;.\;.\;.\;. & $1.589198$\\
Longitude of perihelion & .\;.\;.\;.\;.\;.\;.\;.\;.\;.\;.\;.\;. & $47^\circ\;39'\;36''1$\\
Longitude of the ascending node & .\;.\;.\;.\;.\;.\;.\;.\;. & $92\;53\;44.2$\\
Inclination of the orbit & .\;.\;.\;.\;.\;.\;.\;.\;.\;.\;.\;.\;.\;. & $31\;32\;38.7$\\
\end{tabular}\end{center}
\noindent Direct motion.

\par\indent The constants used to calculate the comet’s coordinates with respect to the equator by Professor \textsc{Gauss}’s method are as follows.

\Needspace{13\baselineskip}
\[
\begin{aligned}
x\cos\tfrac{1}{2}\upsilon^2&=\alpha\sin(\upsilon+138^\circ\;9'\;39''6)\\
y\cos\tfrac{1}{2}\upsilon^2&=\ObservedBeta\sin(\upsilon+59\;53\;53.5)\\
z\cos\tfrac{1}{2}\upsilon^2&=\gamma\sin(\upsilon+355\;26\;5.4)\\
\log\alpha&=0.1319477\\
\log\ObservedBeta&=0.1784770\\
\log\gamma&=9.9865984
\end{aligned}
\]

\Needspace{35\baselineskip}\par\indent In addition to the observations listed above, Professor \textsc{Gauss} received some observations made by Mr. von \textsc{Lindenau} at the Seeberg observatory, by Dr. \textsc{Olbers} in Bremen, and by Mr. \textsc{Burckhardt} in Paris. Burckhardt’s were meridian observations. Mr. \textsc{Nicolai} compared his elements with all these observations and obtained the fine agreement shown below, which leaves hardly anything to be desired.

\begin{center}
\fontsize{10}{12}\selectfont\setlength{\tabcolsep}{5pt}\renewcommand{\arraystretch}{1.15}\begin{tabular}{l|c|c|l}
 & \multicolumn{2}{c|}{Difference} & \\
 & \shortstack{in right\\ascension} & \shortstack{in\\declination} & Observer \\
\hline
1811 Nov. 18 & $-12''4$ & $-138''4$ & v. Zach \\
19 & $+34.7$ & $+19.5$ & v. Zach \\
20 & $-11.9$ & $-5.5$ & v. Zach \\
21 & $-45.0$ & $+4.9$ & v. Zach \\
Dec. 8 & $+9.7$ & .\;.\;.\;. & v. Lindenau \\
9 & $+1.0$ & $-51.4$ & Olbers \\
9 & $+4.9$ & $+22.6$ & v. Lindenau \\
9 & $-5.0$ & $+3.4$ & Gauss \\
11 & $+2.9$ & $-32.7$ & Gauss \\
12 & $+8.5$ & $-27.6$ & Gauss \\
14 & $+23.5$ & $-34.4$ & Olbers \\
14 & $+18.9$ & $-6.7$ & Burckhardt \\
16 & $+28.7$ & $-18.1$ & Olbers \\
22 & $+20.8$ & $+2.4$ & Burckhardt \\
25 & $+25.1$ & $-3.1$ & Burckhardt \\
1812 Jan. 3 & $+2.7$ & $-21.1$ & Olbers \\
3 & $-5.5$ & $+12.9$ & Gauss \\
4 & $+2.8$ & $-12.7$ & Gauss \\
6 & $-35.6$ & $-1.5$ & Olbers \\
\end{tabular}
\end{center}

\par\indent The comet has grown noticeably dimmer since the first observations here, when its brightness was greatest and about $= 0.7$, but on 3 and 4 January it could still be observed just as well as before. It would remain visible for a very long time if there were no reason to fear\SourceLeaf{352} that its steadily fading light will probably hide it from our eyes as early as February. In the meantime, observations will be made much easier by the following ephemeris, which Mr. \textsc{Nicolai} also calculated. According to it, the comet, having already passed right through the Hyades, will move from Taurus into Auriga and will probably disappear from view there.
